\documentclass[10pt,a4paper]{amsart}
\usepackage[margin=19mm]{geometry}
\usepackage{amsmath,amssymb,mathtools}
\usepackage[scr=rsfs,cal=euler]{mathalfa}
\usepackage{xfrac}
\usepackage[unicode,hidelinks,bookmarks=false]{hyperref}
\newcommand{\ie}{i.e.\ }
\newcommand{\cf}{cf.\ }
\newcommand{\resp}{respectively\ }
\newcommand{\Subsection}{\subsection}
\providecommand{\nobreakdash}{\nobreak\hbox{-}}
\setlength{\emergencystretch}{2em}
\multlinegap0pt
\usepackage{amssymb}
\usepackage{mathtools}
\usepackage{xcolor}
\usepackage{enumerate}
\newtheorem{proposition}{Proposition}
\title{Approximate Transitivity of Young Translation on Rough Paths}
\author{Carlo Bellingeri, Paul Gassiat, Hugo Nouaille}
\date{Source: arXiv:2605.30869v1 (CC BY 4.0)}
\begin{document}
\maketitle

\noindent\emph{Source and licence.} Approximate Transitivity of Young Translation on Rough Paths. Version arXiv:2605.30869v1 (CC BY 4.0), distributed by arXiv under the Creative Commons Attribution 4.0 International licence, http://creativecommons.org/licenses/by/4.0/, as recorded in the arXiv licence metadata for this exact version and shown on the abstract page https://arxiv.org/abs/2605.30869v1. Full licence notice: \texttt{licenses/arxiv-2605.30869-CC-BY-4.0.txt}.\par\smallskip

\noindent\textit{Editorial scope.} Counted unit: the article's proposition inversion (source0.tex lines 541--546). The article's complete original proof of it is delivered with it (source0.tex lines 548--599, one \texttt{proof} environment of 52 lines). No mathematics is written, repaired or re-derived: the statement and the proof are the article's own lines, in the article's own notation, and no retained line is substituted or re-flowed.  No further source-local context is needed: every cross-reference of every retained line resolves inside the two delivered spans.  Nothing was omitted from any retained span, so the package carries no omission record.  No retained line contains a citation, so the wrapper carries no bibliography and no bibliography entry.  No retained line contains a figure environment, an image-inclusion command or drawing code, so no image is delivered and the package declares no asset dependency.  Editorial formatting only, in addition: an AMS \texttt{amsart} preamble that restates the article's own declarations and macros used by the retained lines, namely the environments \texttt{proposition} on separate counters, together with the standard packages those lines need.  The retained environments print in this document as Proposition 1 in order of appearance, whereas the article prints the same items with the numbering of its own full document.  The complete native source of the article as distributed in the arXiv e-print for this exact version is bundled unchanged as \texttt{sources/201699/source0.tex}.
\begin{proposition}\label{inversion}
Let $f\colon \mathbb{R}^n \to G^N(\mathbb{R}^d)$ be a $C^1$ function such that
$f(0)=\mathbf{1}$ and $\mathrm{d}^* f_0=\mathrm{d}f_0$ is invertible. Then there exist $R,\varepsilon>0$ such that for every $g \in C^1(\mathbb{R}^n,G^N(\mathbb{R}^d))$ satisfying
$ d_1^R(f,g) \le \varepsilon$,
one has $\mathbf{1}\in g\big(B(0,R/2)\big)$.
\end{proposition}

\begin{proof}
To prove the result we will use the global chart $\log$ to restate the theorem. Setting $F := \log \circ f$ and $G := \log \circ g$, the condition
$\mathbf{1} \in g(B(0,R/2))$ is equivalent to $F(0)=0$ and we search $x^* \in B(0,R/2)$ with
$G(x^*)=0$. Since $\log$ is smooth and for any $h\colon \mathbb{R}^n \to G^N(\mathbb{R}^d)$ one has the explicit formula on $H= \log \circ h$
\[dH_x=\frac{\text{ad}_{H(x)}}{e^{\text{ad}_{H(x)}}- \text{id}}\circ \mathrm{d}^* h_x=\sum_{k=0}^{N-1}\frac{B_k}{k!}\text{ad}^k_{H(x)}\circ \mathrm{d}^* h_x\]
where $\{B_k\}_{k\geq 1}$ are the Bernoulli number with the convention $B_1=- \frac{1}{2}$ and $\text{ad}_{H(x)}$ is the usual adjoint map, the condition $d_1^R(f,g)\leq \varepsilon$ implies
\[
    \sup_{x\in B(0,R)}\|G(x) - F(x)\| +    \sup_{x\in B(0,R)}\|\mathrm{d}G_x - \mathrm{d}F_x\|_{\text{op}} \leq C_{R,N}\,\varepsilon
\]
for some constant $C_{R,N} \geq 1$ depending only on $N$ and $R$. Replacing $\varepsilon$ by $C_{R,N}\varepsilon$ if necessary, we may assume
\begin{equation}\label{eq:closeness}
    \sup_{x\in B(0,R)}\|G(x) - F(x)\| +    \sup_{x\in B(0,R)}\|\mathrm{d}G_x - \mathrm{d}F_x\|_{\text{op}}\leq \varepsilon\,.
\end{equation}
Recalling that $\mathrm{d}F_0= \mathrm{d}^*f_0=\mathrm{d}f_0 $, by continuity of $\mathrm{d}F$ at $0$ we can choose $0<R \leq 1$ small enough  so that
\begin{equation}\label{eq:R}
    \sup_{x \in B(0,R)} \|\mathrm{d}F_x - \mathrm{d}f_0\|_{\text{op}}
    \leq \frac{1}{4\|\mathrm{d}f_0^{-1}\|_{\text{op}}}\,,
\end{equation}
which is possible by continuity. Set $\varepsilon := R/(4\|(\mathrm{d}f_0)^{-1}\|)$. We claim that the map
\[
    T(x) := x - ((\mathrm{d}f_0)^{-1})G(x)
\]
is a contraction on $B(0,R)$ mapping into $B(0,R/2)\subset B(0,R)$, so that its unique fixed point yields the desired zero of $G$.  Indeed, for
$x\in B(0,R)$, one has 

\[T(x) = \mathrm{d}f_0^{-1}(\mathrm{d}f_0\,x - G(x))= \mathrm{d}f_0^{-1}\left((\mathrm{d}f_0\,x - F(x)) + (F(x)-G(x))\right)\]

By the mean value inequality applied to $x\mapsto F(x)-\mathrm{d}f_0\,x$, which vanishes at $0$ together with  the property \eqref{eq:R}, one has for all $x\in B(0,R)$
\begin{equation}\label{eq:MVT}
    \|F(x) - \mathrm{d}f_0\,x\| \leq \frac{\|x\|}{4\|(\mathrm{d}f_0)^{-1}\|_{\text{op}}}\,.
\end{equation}
Then from estimates \eqref{eq:MVT} and \eqref{eq:R} we obtain
\[
    \|T(x)\| \leq \|(\mathrm{d}f_0)^{-1}\|_{\text{op}}\!\left(\frac{R}{4\|(\mathrm{d}f_0)^{-1}\|_{\text{op}}}
    + \varepsilon\right) = \frac{R}{2},
\]
For the contraction property, take  $x,y\in B(0,R) $ and $z\in B(0,R)$. From  \eqref{eq:R} and \eqref{eq:closeness} one has
\[
    \|\mathrm{d}f_0 - \mathrm{d}G_z\|
    \leq \|\mathrm{d}f_0 - \mathrm{d}F_z\| + \|\mathrm{d}F_z-\mathrm{d}G_z\|
    \leq \frac{1+R}{4\|(\mathrm{d}f_0)^{-1}\|_{\text{op}}},
\]
and therefore, writing \[T(x)-T(y) = (\mathrm{d}f_0)^{-1}\int_0^1(\mathrm{d}f_0 -
\mathrm{d}G_{y+t(x-y)})(x-y)\,\mathrm{d}t\,,\]
we obtain
\[
    \|T(x)-T(y)\| \leq \frac{1+R}{4}\,\|x-y\| \leq \frac{1}{2}\,\|x-y\|,
\]
where the last inequality uses $R\leq 1$. By the Banach fixed point theorem, $T$ has a
unique fixed point $x^*\in B(0,R) $, which by the self-mapping property satisfies $x^*\in B(0,R/2)$ and $G(x^*)=0$, i.e.\ $g(x^*)=\mathbf{1}$.
Hence $\mathbf{1}\in g(B(0,R/2))$.
\end{proof}

\end{document}
